\section{Quantum connections and one stabilization}
\label{sec:atomic-one-step}

We first define the two rank-two counts, then establish their intrinsic
blowup and projective-bundle formulas.  Low-dimensional vanishing and one
explicit residue calculation give the cubic obstruction.  The abstract
additive formulation follows that proof.

\subsection{Whole primary connection blocks}

Let \(Y\) be smooth and projective.  We first make three setup choices.
Pass to numerical Novikov variables, set the odd bulk variables to zero, and
restrict the underlying cohomology space to \(H^{\mathrm{ev}}(Y)\).  The
last restriction is essential: restricting the bulk base alone would still
leave odd cohomology in the fibre.  Use graded completions as in
\cite[Section~2.2]{IritaniBlowup}, and let \(K_Y\) be the fraction field of
the numerical divisor-equation-reduced coefficient domain described below.
Fix an algebraic closure \(\overline K_Y\).  We form
all spectral blocks only after base change to \(\overline K_Y\), so no later
scalar extension can split a factor that has already been counted as one
block.

For a homogeneous even cohomology basis \((\phi_\alpha)\), our conventions are
\[
 \nabla_{\partial_{t^\alpha}}
 =\partial_{t^\alpha}+z^{-1}(\phi_\alpha\star_t),\qquad
 \nabla_{z\partial_z}=z\partial_z-z^{-1}(E_Y\star_t)+\mu_Y,
\]
where \(n=\dim Y\), \(\mu_Y|_{H^k(Y)}=(k-n)/2\), and
\[
 E_Y=c_1(Y)+\sum_\alpha
 \left(1-\frac{\deg\phi_\alpha}{2}\right)t^\alpha\phi_\alpha.
\]
Horizontal sections therefore satisfy
\(z^2\partial_z y=((E_Y\star_t)-z\mu_Y)y\).
Let \(U\) denote multiplication by the Euler field on \(H^{\mathrm{ev}}(Y)\).
For a primary block with eigenvalue \(\lambda\), centering removes the
scalar exponential factor \(e^{-\lambda/z}\); its leading operator becomes
\(N=U-\lambda I\).

After base change to \(\overline K_Y\), each generalized eigenspace of \(U\) is
a summand of the full formal connection, not just an invariant subspace of
\(U\):

\begin{proposition}[Formal spectral connection splitting]
\label{prop:generic-spectral-connection-splitting}
\coverage{fragment}%
\lean{genericSpectralConnection_normalizedSplitting}%
Suppose the leading operator of a finite free formal QDM has block-diagonal
form \(U_i=\lambda_iI+N_i\), where each \(N_i\) is nilpotent and
\(\lambda_i-\lambda_j\) is a unit for \(i\ne j\).  There is a unique gauge
\(G=I+O(z)\), regular in \(z\) and block off-diagonal in positive orders,
which makes the \(z\)-connection block diagonal.  In this gauge every bulk
and Novikov direction is block diagonal as well.  The horizontal pairing is
block diagonal and nondegenerate on each factor.
\end{proposition}

\begin{proof}
At order \(z^{n-1}\), the \((i,j)\)-entry of the gauge equation is a Sylvester
equation
\[
 U_iX-XU_j=B_{n,ij}.
\]
Its operator is \((\lambda_i-\lambda_j)I\) plus the nilpotent operator
\(X\mapsto N_iX-XN_j\), hence is invertible.  Recursion gives the normalized
gauge, and the same equation with zero right-hand side proves uniqueness.

Choose frames of the generalized-eigenspace subbundles after the splitting
extension.  If an off-diagonal block of a base-direction connection had a
first nonzero Laurent coefficient \(Xz^m\), the lowest off-diagonal coefficient
of flatness would be \(U_iX-XU_j=0\), contrary to Sylvester invertibility.
Thus every base direction preserves the same factors.  Applied coefficientwise
to pairing horizontality, the identical argument kills the off-diagonal pairing
blocks.  Nondegeneracy of the total pairing then makes every diagonal
restriction nondegenerate.
\end{proof}

We call these factors the \emph{generic even QDM blocks} of \(Y\).
Each is a whole primary spectral summand; individual Jordan blocks within
one primary summand are not counted separately.  A block
retains its connection, pairing, and centered Euler endomorphism, but not an
absolute cohomological grading.  A block type may retain additional marked
data.  If it retains grading, a morphism may carry a recorded integral
suspension, and the invariant must be invariant under that suspension.  The block
type used below forgets absolute grading.  Its morphisms preserve the
retained data and are regular connection isomorphisms after scalar extension
and formal change of bulk coordinates.  In particular, their gauges use
neither negative nor fractional powers of the connection variable \(z\).

\subsection{Rank-two residues and two counts}

Let \(E\) be a generic even block with centered connection
\begin{equation}
 z\partial_z y=A(z)y,
 \qquad A(z)=z^{-1}N+A_0+zA_1+z^2A_2+\cdots.
 \label{eq:atomic-block-expansion}
\end{equation}
Assume \(\rk E=2\), \(N\ne0\), and \(N^2=0\).  Then
\(L=\im N=\ker N\) is a line.  The following matrix calculation explains
the modification.  In a frame adapted to \(L\), suppose that \(A_0\)
preserves this line, so that
\[
 N=\begin{pmatrix}0&\nu\\0&0\end{pmatrix},\qquad
 A_0=\begin{pmatrix}a&b\\0&d\end{pmatrix},\qquad
 (A_1)_{21}=c.
\]
Changing the frame by \(S=\diag(1,z)\) removes the pole and gives residue
\[
 R=\begin{pmatrix}a&\nu\\c&d-1\end{pmatrix}.
\]
The lower-left entry explains why the first-order coefficient \(A_1\)
is needed; the subtraction of one comes from differentiating the new
frame.  Eigenvalues \(0,1\) represent the same class modulo \(\Z\)
and do not qualify, whereas \(0,2/3\) represent distinct classes.

The canonical elementary modification along
this line is
\begin{equation}
 E^\sharp=\{s\in E:s\bmod z\in L\}.
 \label{eq:atomic-modification}
\end{equation}

For this modification to remove the pole, the regular coefficient \(A_0\) must
preserve \(L\).  Pairing horizontality supplies exactly that condition.

\begin{lemma}[The regular coefficient preserves the nilpotent line]
\label{lem:A0preserve}
\coverage{conditional_deduction}%
\lean{atomicRankTwo_regularCoefficient_preserves_nilpotentLine,%
  atomicRankTwo_pairingEquations_of_loopHorizontality}%
Horizontality of the Poincar\'e pairing implies \(A_0L\subset L\).
\end{lemma}

\begin{proof}
Separated spectral factors are orthogonal for the horizontal Poincar\'e
pairing: coefficientwise horizontality kills every off-diagonal block by an
invertible Sylvester operator.  Hence the restriction to \(E\) is
nondegenerate.  Write it as \(P(z)=P_0+zP_1+\cdots\).  Horizontality gives
\[
 z\partial_zP(z)+A(z)^TP(z)+P(z)A(-z)=0.
\]
The \(z^{-1}\) coefficient gives \(N^TP_0=P_0N\).  Thus \(L\) is isotropic
and, in dimension two, \(L^\perp=L\).  The constant coefficient is
\[
 A_0^TP_0+P_0A_0+N^TP_1-P_1N=0.
\]
Evaluating this identity on \((x,x)\), where \(x\in L=\ker N\), gives
\((A_0x,x)=0\).  Hence \(A_0x\in L^\perp=L\).
\end{proof}

Choose a frame with \(L=\langle e_1\rangle\) and
\(Ne_2=\nu e_1\), \(\nu\ne0\).  The modified lattice has frame
\(e_1,ze_2\), so
\begin{equation}
 A^\sharp=S^{-1}AS-S^{-1}z\partial_zS,
 \qquad S=\diag(1,z).
 \label{eq:Asharp}
\end{equation}
The term \(z^{-1}N\) becomes regular.  The only possible new pole is the
\(E_{21}\)-coefficient of \(A_0\), which vanishes by
Lemma~\ref{lem:A0preserve}.  Thus
\(A^\sharp=R+zR_1+\cdots\).  Define
\begin{equation}
 \delta^\sharp(E):=(\tr R)^2-4\det R.
 \label{eq:res-disc}
\end{equation}

\begin{lemma}[Persistence of a cyclic double-primary cluster]
\label{lem:cyclic-primary-persistence}
\coverage{absent}%
\uses{prop:generic-spectral-connection-splitting}%
Over a characteristic-zero formal bulk germ, let a rank-two spectral
cluster of a flat QDM be separated from its complement at the closed point.
Suppose its centered leading operator specializes to a nonzero square-zero
operator, and its scalar-centered base equations have at most a simple
\(z\)-pole.  The cluster remains cyclic and its centered leading operator
remains nonzero square-zero throughout the germ.
\end{lemma}

\begin{proof}
Write a centered base equation as
\(\partial y=(z^{-1}C_\partial+C_{\partial,0}+O(z))y\).  The
\(z^{-2}\) coefficient of flatness gives \([N,C_\partial]=0\), and the trace
of the \(z^{-1}\) coefficient gives \(\tr C_\partial=0\); here the base
equation has been scalar-centered.  Nilpotence of \(N\) away from the
base point is not yet known, so the commutant is computed over the
complete local ring \(\mathcal O\) of the formal bulk germ.  The rank-two
spectral cluster continues over \(\mathcal O\) because it is separated
from the complementary spectral clusters at the base point; \(N\) is the centered
\(z^{-1}\) coefficient of that continuation, so \(\tr N=0\), and it
specializes at the closed point to a nonzero rank-one nilpotent.  Choose
\(v\) with \(v,Nv\) a basis modulo the maximal ideal; by Nakayama it is a
basis over \(\mathcal O\), so \(N\) is cyclic.  Any \(C\) commuting with a
cyclic \(N\) is determined by \(Cv=\alpha v+\beta Nv\) and equals
\(\alpha I+\beta N\), whether or not \(N\) is nilpotent.  Hence
\([N,C_\partial]=0\) and \(\tr C_\partial=0=\tr N\) give
\(C_\partial=q_\partial N\) with \(q_\partial\in\mathcal O\).  The same
\(z^{-1}\) coefficient gives
\[
 \partial N=-q_\partial N+
 [N,q_\partial A_0-C_{\partial,0}].
\]
Writing \(B=q_\partial A_0-C_{\partial,0}\), one has
\(\partial(N^2)=-2q_\partial N^2+[N^2,B]\); uniqueness for the formal
coefficient recursion preserves \(N^2=0\).  An entry nonzero at the closed
point remains a unit, so \(N\) has rank one throughout the formal germ.
\uses{prop:generic-spectral-connection-splitting}%
\end{proof}

\begin{proposition}[Invariance of the modified residue]
\label{prop:rank2-rigidity}
\coverage{conditional_deduction}%
\lean{atomicRankTwo_modifiedBase_regular_of_flat_connection,%
  atomicRankTwo_modifiedResidue_lax_of_flat_connection,%
  atomicRankTwo_residueDiscriminant_frozen_by_horizontality_and_flatness,%
  atomicRankTwo_residueDiscriminant_constant_over_formal_germ,%
  atomicRankTwo_modifiedBase_pole_eq_zero,%
  atomicRankTwo_residueDiscriminant_constant_along_base}%
Under formal change of even bulk coordinates, the modified connection is
regular in every base direction and
\begin{equation}
 \partial R=[G_\partial,R]
 \label{eq:atomic-lax}
\end{equation}
for a regular matrix \(G_\partial\).  Consequently
\(\partial(\delta^\sharp)=0\).  Regular connection isomorphisms preserve the
domain and value of \(\delta^\sharp\).
\end{proposition}

\begin{proof}
Lemma~\ref{lem:cyclic-primary-persistence} preserves the nilpotent line.
Its centralizer argument writes the centered leading base coefficient as
\(q_\partial N\).  In a frame with \(N=\nu E_{12}\), the modification
\(\diag(1,z)\) makes this upper-right pole regular.  Only the lower-left
entry of the regular base coefficient can then create a pole, of the form
\(z^{-1}kE_{21}\).  If
\(R=\left(\begin{smallmatrix}a&\nu\\c&d\end{smallmatrix}\right)\), the
\(z^{-1}\) flatness equation is
\(kE_{21}+[R,kE_{21}]=0\).  Its diagonal entries are \(\nu k\) and
\(-\nu k\), so \(k=0\).  The constant term of the flatness equation is now
\eqref{eq:atomic-lax}; trace and determinant are constant under infinitesimal
conjugation.  A regular isomorphism carries \(N\), \(L\), and the modified
lattice to their counterparts and conjugates the residue.  Indeed, if
\(G,G^{-1}\) are regular on the original lattices, the induced map
\(S_{\rm target}^{-1}GS_{\rm source}\) and its inverse are regular on the
modified lattices.  Reduction modulo \(z\) gives the conjugating matrix;
its lower-left entry can involve the first coefficient of \(G\), so it
need not be \(G(0)\) in the original frames.
\uses{lem:cyclic-primary-persistence,lem:A0preserve}%
\end{proof}

\begin{proposition}[Exponent interpretation]
\label{prop:residue-discriminant-exponents}
\coverage{fragment}%
\lean{residueDiscriminant_eq_squared_eigenvalue_separation}%
\imports{PutSinger:ch-3}%
If \(\rho_1,\rho_2\) are the eigenvalues of \(R\), their classes modulo
\(\Z\) are the formal exponent classes of the original centered block and
\[
 \delta^\sharp=(\rho_1-\rho_2)^2.
\]
\end{proposition}

\begin{proof}
The elementary modification uses integral powers of \(z\), so it preserves
formal exponent classes modulo \(\Z\).  The modified connection is regular
singular with residue \(R\); the regular-singular classification identifies
the exponent classes with the residue eigenvalues modulo \(\Z\)
\cite[Chapter~3]{PutSinger}.  The displayed identity is the quadratic
discriminant formula.
\end{proof}

For such a block \(E\), write \(\rho_1,\rho_2\) for the eigenvalues of its
modified residue in the
fixed algebraic closure \(\overline K\) of the generic coefficient field.
Define the invariant on each block by
\begin{equation}
 w_{\mathrm{exp}}(E)=
 \begin{cases}
 1,&\rk E=2,\ N\ne0,\ N^2=0,\
      [\rho_1]\ne[\rho_2]\text{ in }\overline K/\Z,\\
 0,&\text{otherwise},
 \end{cases}
 \label{eq:atomic-fold}
\end{equation}
and put \(I_{\mathrm{exp}}(Y)=\sum_E w_{\mathrm{exp}}(E)\).
Call a block \emph{eligible} if it has rank two and \(N\ne0\), \(N^2=0\).
Define the \emph{lattice count} by
\[
 I_{\mathrm{lat}}(Y)
 =\#\{E\text{ eligible}:\delta^\sharp(E)\ne0\}.
\]
Nonzero discriminant means that the two eigenvalues of the canonical
modified residue are distinct.  The stronger test
\(\delta^\sharp\notin\{n^2:n\in\Z\}\) means that their classes modulo
\(\Z\) are distinct.  Thus \(I_{\mathrm{exp}}\leq I_{\mathrm{lat}}\).
Both counts are preserved by regular lattice isomorphisms and by common
scalar shifts of the residue.  For \(I_{\mathrm{lat}}\), preservation of
the canonical modified lattice is essential: separate integral shifts of
the two residue eigenvalues need not preserve their difference.
Neither cyclic persistence nor residue conjugacy uses nonresonance.

\subsection{Compatibility with blowups and projective bundles}

To use the cited decomposition theorems, we must check that they preserve all
the data seen by the invariant.  Four issues arise: coefficient fields,
restriction to even cohomology, regularity in \(z\), and faithful embedding of
the center coefficients.
The last issue is the delicate one: Iritani's center Novikov map may have a
kernel, so the independent target coordinates must recover the missing
numerical curve-class data.
\paragraph{Coefficient fields.}
The intrinsic ample-adic completion and the Laurent neighborhood of a
fibre or exceptional cusp need not map to one another.  Iritani's
display~(1.1) supplies the required common coefficient maps
\cite[(1.1)]{IritaniBlowup}; injectivity for each center summand is
established below.  For projective bundles the common source is
\(\C[z,q][[Q,\widehat\tau]]\).  On the blowup side the ambient term uses the
corresponding graded source.  A center term begins with its intrinsic reduced
coefficient ring from Remark~2.3, maps to the possibly smaller image ring of
Remark~5.6, and is then pulled to a separate target ring for that summand in
the external direct sum.  Embed the generic fields of the ambient common image
and of these separate center images in a fixed algebraically closed overfield
and form every primary block there.  Base
change after this splitting extension preserves the rank, nilpotent Jordan
type, regular formal gauge, and formal exponent classes of each block; forming
primary factors over a nonsplitting field before extension would not.

\paragraph{Restriction to even cohomology.}
The comparison is restricted to even cohomology, not merely to the
even bulk base.  The supermodule conventions and the pull-push and Fourier
formulas in Iritani's Section~5.8 and Iritani--Koto's construction are
parity-preserving.  After the odd bulk variables are set to zero, the
comparison and its inverse therefore restrict to mutually inverse maps on
\(H^{\mathrm{ev}}\).  This is the QDM on which the additive invariant is defined.

\paragraph{Regularity in \(z\).}
The comparisons are regular in \(z\).  Iritani's Theorem~5.18 and
Iritani--Koto's Theorem~5.1 are respectively defined over rings of the form
\[
 \C[z]((q_{\mathrm{exc}}^{-1/\mathfrak s}))[[Q,\widetilde\tau]],
 \qquad
 \C[z]((q_{\mathrm{fib}}^{-1/r'}))[[Q,\widehat\tau]],
\]
where Iritani--Koto take \(r'=r\) when \(r-1\) is even and \(r'=2r\) when
\(r-1\) is odd.  Here \(\mathfrak s\in\{c-1,2(c-1)\}\) is Iritani's
parity-dependent ramification index from \cite[(5.11)]{IritaniBlowup}.  Their
inverses use no negative powers of \(z\); see also
\cite[Remark~5.3]{IritaniKoto}.  The maps are homogeneous but need not have
degree zero: the projective-bundle map has degree \(-(r-1)\), while Iritani's
ambient and center components have degrees \(0\) and \(-c\), respectively
\cite[Theorem~5.1(3)]{IritaniKoto}\cite[Theorem~5.18(3)]{IritaniBlowup}.
The cited maps intertwine the actual \(z\)-connections.  A common grading
normalization can translate the residue by a scalar, preserving its
discriminant; no eigenline is independently regraded.  The graded completed rings embed coefficientwise
in \(\overline K[[z]]\); no identification with an ordinary tensor-product
completion is needed.

\paragraph{Faithful center coordinates.}
Iritani's center map need not be injective
\cite[(5.15)]{IritaniBlowup}.  After numerical reduction it has the form
\begin{equation}
 Q_Z^d\longmapsto
 Q^{i_*d}q_{\mathrm{exc}}^{-\rho_Z\cdot d/(c-1)}.
 \label{eq:center-novikov-map-atomic}
\end{equation}
Here \(\rho_Z=c_1(N_{Z/Y})\).
The raw map is not the coefficient map used below.  Remark~2.3 gives the
intrinsic divisor-equation-reduced coefficient ring, whose Novikov generators
are the combined symbols
\[
 X_d=Q_Z^d\exp\!\bigl(\sigma^{(2)}\cdot d\bigr);
\]
the divisor coordinates \(\sigma^{(2)}\) are not separate elements of this
ring.  Remark~5.6 places its pushforward under the raw Novikov map over the
image of that reduced ring \cite[Remarks~2.3 and~5.6]{IritaniBlowup}.  In the
external direct sum of Section~5.8.2, by contrast, the variables
\(s_j\in H^*(Z)\) are independent coordinates on the common comparison target,
and the center parameter is identified as
\(\varsigma_j=\varsigma_j^\circ+s_j\)
\cite[(5.47)--(5.48)]{IritaniBlowup}.  The divisor components of these
independent target coordinates retain the curve-class information lost by
\eqref{eq:center-novikov-map-atomic}.

The distinction is modeled by \(x,y\mapsto q\), which identifies two
generators, and \(x\mapsto qe^s,\ y\mapsto qe^t\), which retains independent
characters.  Source coefficients must not already depend on the target
variables \(s,t\).  The comparison uses the following conventions:
\begin{center}
\small
\begin{tabular}{@{}p{.22\linewidth}p{.72\linewidth}@{}}
\toprule
Object & Required data \\
\midrule
Source & Reduced graded coefficients \(\C[[X_d,\sigma']]_{\rm gr}[\sigma^0]\);
no independent source divisor variables. \\
Target & Independent ambient and center-copy coordinates \(t,s_0,\ldots,s_{c-2}\);
all center divisor directions retained. \\
Scalars & The displayed exceptional or fibre Novikov Laurent extensions and
ramification; their \(z\)-free fraction fields in a common algebraic closure. \\
Connections & Graded completed \(z\)-enhancements embedded coefficientwise
in the common field's \(\overline K[[z]]\); comparison and inverse regular. \\
Derivations & Bulk partials, numerical Novikov logarithmic derivatives and
the actual \(z\)-connection, with coordinate-change chain rules. \\
Grading & Only common scalar residue shifts; no independent shifts of residue
eigenlines. \\
\bottomrule
\end{tabular}
\end{center}

\begin{lemma}[Faithful center base change]
\label{lem:faithful-center-base-change}
\coverage{fragment}%
\lean{faithfulCenterBaseChange_targetOnly_injective}%
Let \(R_{Z,\mathrm{src}}^{\mathrm{red}}
=\C[[X_d,\sigma']]_{\mathrm{gr}}[\sigma^0]\) be the coefficient domain
whose graded completed \(z\)-enhancement is Iritani's reduced source ring
\[\C[z][[Q_Ze^{\sigma^{(2)}},\sigma']][\sigma^0].\]
Let \(A_j\) be the
numerical coefficient domain of the external direct sum in Section~5.8.2,
with the independent target coordinates \(t,s_0,\ldots,s_{c-2}\) and with
\(z\) omitted.  The pullback determined by the center Novikov map and
\(\varsigma_j=\varsigma_j^\circ+s_j\) gives a continuous homomorphism
\[
 \iota_j:R_{Z,\mathrm{src}}^{\mathrm{red}}\longrightarrow A_j.
\]
On a combined divisor-equation generator it is
\begin{equation}
 X_d=Q_Z^d\exp\!\bigl(\sigma^{(2)}\cdot d\bigr)
 \longmapsto
 Q^{i_*d}q_{\mathrm{exc}}^{-\rho_Z\cdot d/(c-1)}
 \exp\!\bigl((\varsigma_j^{\circ,(2)}+s_j^{(2)})\cdot d\bigr).
 \label{eq:faithful-center-map}
\end{equation}
This map is injective.  In particular, it induces a field embedding
\[
 \operatorname{Frac}(R_{Z,\mathrm{src}}^{\mathrm{red}})
 \lhook\joinrel\longrightarrow
 \operatorname{Frac}(A_j).
\]
After this field extension, each center summand in Iritani's comparison is the
faithful scalar extension, followed by the formal coordinate change
\(\varsigma_j\), of the intrinsic generic numerical even QDM of \(Z\).
\end{lemma}

\begin{proof}
All completions here use the graded convention of
\cite[Section~2.2]{IritaniBlowup}.  The pullback is defined on the reduced
ring by the String and Divisor Equations, as explained after (5.36) and
(5.47) there.  It need not be defined on the unrestricted bulk power-series
ring.  We prove injectivity using the first nonzero Novikov degree, retaining
the full exponentials in the target divisor coordinates.
Choose integral divisors \(D_1,\ldots,D_\rho\) whose pairings separate
\(N_1(Z)\), and use their dual coordinates in \(s_j^{(2)}\).  The variables
\(s_j\) occur only in the target \(A_j\), not as coefficients in
\(R_{Z,\mathrm{src}}^{\mathrm{red}}\).  This asymmetry is essential.

Filter by an integral ample degree pulled back from \(Y\).  Every fibre of
the untagged monomial map in a bounded piece is finite: the restriction
of an ample divisor from \(Y\) is
ample on \(Z\), and its bounded slice of \(\overline{NE}(Z)\) contains only
finitely many lattice points.  This filtration is respected by the target
map because \(H\cdot i_*d=(i^*H)\cdot d\), and the fixed shifts have
nonnegative Novikov degree.

For a fixed curve class, the coefficient of an element of the reduced ring
is polynomial in the nondivisor and unit coordinates.  Indeed, an element
has finitely many homogeneous degrees and bounded unit degree; every other
even nondivisor coordinate has strictly negative degree.  Thus only
finitely many of its monomials occur at that curve class.  This is where
the reduced, graded source is essential.

Suppose a nonzero source element maps to zero and take its first nonzero
Novikov-degree part.  At that degree only the Novikov-degree-zero part of
\(\varsigma_j^\circ\) enters.  Translation by its nondivisor and unit
components is an injective substitution on the polynomial coefficients,
over the target Laurent coefficient field.  Its divisor component
multiplies each exponential by a nonzero scalar.  Group the terms by their
common untagged target monomial.  On one nonzero group the
remaining relation has the form
\[
 \sum_{d\in F}a_d
 \exp\!\left(\sum_{k=1}^\rho(D_k\cdot d)s_{j,k}\right)=0,
 \qquad a_d\ne0,
\]
over a characteristic-zero coefficient domain, including the remaining bulk
coordinates, that does not contain the variables \(s_{j,k}\).
The exponent vectors are distinct.  Choose an integral
one-parameter restriction on which they remain distinct.  Derivatives of
orders \(0,\ldots,|F|-1\) at the origin give a Vandermonde matrix, hence every
\(a_d\) is zero, a contradiction.  This proves injectivity on the full
reduced source.  We have not truncated the divisor exponentials by bulk
order; doing so could lose the distinction between curve classes.

Remark~2.3 of \cite{IritaniBlowup} gives the intrinsic reduced source, and
Remark~5.6 gives its image under the raw Novikov map~(5.15).
Although that map can be noninjective on the unrestricted source, it retains
the independent \(\sigma^{(2)}\) coordinates in its target.  Its restriction
to the reduced source is injective by the same argument, with zero shift.
Equations~(5.47)--(5.48) then pull that image into the external direct sum over
\(A_j\), where the variables \(s_j\) are independent, and identify the direct
sum with the blowup module.  Associativity of scalar extension therefore
identifies the center summand with scalar extension directly along the
injective composite \(\iota_j\).  Passing to fraction fields is faithfully
flat.  Finally, the inverse of the combined formal coordinate change
\((\tau,\varsigma_0,\ldots,\varsigma_{c-2})\), supplied by
Theorem~5.18(7) and Section~5.8.2, transports the same statement to the
original blowup coordinates.
\end{proof}

\begin{proposition}[Intrinsic formulas for the two counts]
\label{prop:intrinsic-count-formulas}
\coverage{absent}%
\uses{lem:faithful-center-base-change,prop:rank2-rigidity}%
\imports{IritaniKoto:thm-5.1,IritaniBlowup:thm-5.18}%
For either \(I=I_{\mathrm{exp}}\) or \(I=I_{\mathrm{lat}}\), a smooth
codimension-\(c\) center \(Z\subset Y\) and a rank-\(r\) vector bundle \(V\)
satisfy
\begin{equation}
 I(\Bl_ZY)=I(Y)+(c-1)I(Z),\qquad
 I(\PP_Y(V))=rI(Y).
 \label{eq:intrinsic-count-formulas}
\end{equation}
The formulas are componentwise for disconnected centers.
\end{proposition}

\begin{proof}
\emph{Projective bundles.}
For a rank-\(r\) bundle, Iritani--Koto's Theorem~5.1 gives \(r\) copies of
the base QDM after Laurent extension \cite[Theorem~5.1]{IritaniKoto}.  Its
invertible joint bulk Jacobian gives independent unit coordinates on the
\(r\) target factors.  In the \(j\)-th factor, changing its unit coordinate
adds an independent scalar \(u_jI\) to Euler multiplication.  For two finite
spectra, the resultant of their translated characteristic polynomials is a
nonzero polynomial in \(u_i-u_j\), with leading coefficient one.  Hence the
spectra of distinct factors are disjoint at the generic point: their primary
blocks do not fuse when the direct sum is formed.  Regularity, restriction to
even cohomology, the fixed splitting field, and invariance under common scalar grading shifts therefore
carry the invariant blockwise across the comparison, and additivity gives
\eqref{eq:intrinsic-count-formulas}.  Tensoring the bundle by a line bundle, when needed
for the global-generation convention, changes neither the projective bundle
nor the intrinsic Novikov embedding
\cite[Remark~1.2 and Remark~5.2]{IritaniKoto}.

\smallskip
\noindent\emph{Blowups.}
For a codimension-\(c\) blowup, Iritani's Theorem~5.18 gives one ambient QDM
and \(c-1\) center QDMs \cite[Theorem~5.18]{IritaniBlowup}.  Parts (4), (6),
and (7) give the comparison asymptotics, the leading reconstruction
parameters, and the invertible combined bulk Jacobian, respectively.  The
independent unit coordinates supplied by part~(7) give the same nonzero
resultant argument, so the spectra of the ambient and individual center
summands are pairwise disjoint generically.  Regularity, restriction to even
cohomology, the fixed splitting field, invariance under common scalar grading shifts, and
Lemma~\ref{lem:faithful-center-base-change} for each center copy then give
\eqref{eq:intrinsic-count-formulas}.  Faithful center base change identifies each of the \(c-1\) contributions
with \(I(Z)\).  Each comparison is evaluated over its own common
coefficient ring; the proof never composes Laurent comparisons recursively.
\uses{lem:faithful-center-base-change,prop:rank2-rigidity}%
\imports{IritaniKoto:thm-5.1,IritaniBlowup:thm-5.18}%
\end{proof}

\subsection{A rational rank-two calculation}

The zero-primary block can be separated without splitting the complementary
eigenlines.  This gives one calculation for three Fano families.

\begin{proposition}[A universal modified residue]
\label{prop:universal-rank-two-residue}
\coverage{absent}%
\uses{lem:A0preserve}%
For \(s=2a+b\), \(sq\ne0\), consider
\[
 z^2\partial_z y=(U_{a,b}+zD)y,\qquad
 U_{a,b}=\begin{pmatrix}
 0&aq&0&a^2q^2\\1&0&bq&0\\0&1&0&aq\\0&0&1&0
 \end{pmatrix},\qquad D=\tfrac12\diag(3,1,-1,-3).
\]
Its characteristic polynomial is \(T^2(T^2-sq)\).  Its zero-primary block
has rank two, nonzero square-zero leading term and canonical modified residue
\begin{equation}
 R_{a,b}=
 \begin{pmatrix}
 -(2a+3b)/(2s)&1\\
 -4a^2/s^2&(b-2a)/(2s)
 \end{pmatrix},\qquad
 \delta^\sharp(a,b)=\frac{4(b-2a)}{b+2a}.
 \label{eq:universal-residue}
\end{equation}
Here \(a,b,q\) are constant under \(\partial_z\).
\end{proposition}

\begin{proof}
The columns of
\[
 C=\begin{pmatrix}
 aq&0&0&-(a+b)q\\
 0&(a+b)q&-aq&0\\
 1&0&0&1\\0&1&1&0
 \end{pmatrix},\qquad \det C=-s^2q^2,
\]
give \(\im U^2\oplus\ker U^2\).  Direct multiplication gives
\[
 J=C^{-1}UC=J_+\oplus N,\qquad
 J_+=\begin{pmatrix}0&sq\\1&0\end{pmatrix},\quad
 N=\begin{pmatrix}0&1\\0&0\end{pmatrix}.
\]
Put \(B=C^{-1}DC\).  The normalized gauge \(G=I+zX+z^2Y+\cdots\),
with zero diagonal blocks in positive orders, starts with
\[
 X_{+0}=\diag(2a/s,2/(qs)),\qquad
 X_{0+}=\diag(-2/(qs),-2a/s).
\]
These matrices solve \(J_iX_{ij}-X_{ij}J_j=-B_{ij}\).
Writing the transformed coefficient as \(J+zB_1+z^2B_2+\cdots\), its
gauge equation gives
\[
 B_1=B+[J,X],\qquad B_2=BX-XB_1-X+[J,Y].
\]
The derivative contributes \(-X\).  Its diagonal blocks vanish, as do
those of \(XB_1\) and \([J,Y]\), so
\[
 (B_1)_{00}=\diag\left(-\frac{2a+3b}{2s},\frac{2a+3b}{2s}\right),
 \qquad ((B_2)_{00})_{21}=(B_{0+}X_{+0})_{21}
 =-\frac{4a^2}{s^2}.
\]
Modification by \(\diag(1,z)\) gives \(R_{a,b}\).  Its characteristic
polynomial is \(T^2+T+(10a-3b)/(4s)\), proving the discriminant formula.
For an original \(z^2F\) term, add \((C^{-1}FC)_{ii}\) to \((B_2)_{ii}\).
Thus only the first off-diagonal gauge coefficient is needed.
\end{proof}

\subsection{The cubic block}
\label{subsec:shared-cubic-packet}

Only a rank-two block can contribute to \(I_{\mathrm{exp}}\).  For a cubic
threefold, the calculation below isolates the unique such candidate and then
computes the residue of its elementary modification.

Let \(X\subset\PP^4\) be a smooth cubic threefold, let \(P=H|_X\), and put
\(q=Q^\ell\) for the line class.  Beauville's formulas
give \cite[main theorem and (2.1)--(2.3)]{Beauville}
\begin{equation}
 P\star P=P^2+6q,\qquad
 P\star P^2=P^3+15qP,\qquad
 P\star P^3=6qP^2+36q^2.
 \label{eq:beauville-cubic-products}
\end{equation}
Thus, in the basis \((1,P,P^2,P^3)\),
\begin{equation}
 U_X=2\begin{pmatrix}
 0&6q&0&36q^2\\1&0&15q&0\\0&1&0&6q\\0&0&1&0
 \end{pmatrix},
 \qquad D_X=\frac12\diag(3,1,-1,-3),
 \label{eq:cubic-small-even-connection}
\end{equation}
and \(z^2\partial_zy=(U_X+zD_X)y\).  This agrees with
\cite[Section~3]{Cai}.

The change of basis \(\diag(1,2,4,8)\) changes \(U_X=2U_{6,15}\)
to \(U_{24,60}\) and leaves \(D_X\) unchanged.  The universal calculation
therefore applies.  Conjugating its residue by \(\diag(2,1)\)
gives the original cubic normalization:
\begin{equation}
 R_X
 =\begin{pmatrix}-19/18&2\\-8/81&1/18\end{pmatrix}.
 \label{eq:RX}
\end{equation}
Its characteristic polynomial is
\((T+1/6)(T+5/6)\).  Thus the two exponent classes differ by
\(2/3\notin\Z\), so the block is counted.  Equivalently,
\begin{equation}
 \boxed{\delta^\sharp(X)=\frac49},
 \qquad I_{\mathrm{exp}}(X)=1.
 \label{eq:cubic-delta}
\end{equation}

\paragraph{From the small connection to the generic block.}
After inverting \(q\), at the small quantum point the rank-two cluster is separated from the
other eigenvalues.  The cyclic-centralizer argument in
Lemma~\ref{lem:cyclic-primary-persistence} continues that cluster over the formal
bulk ring and proves that its centered leading term remains nonzero and
square-zero.  The modified residue is then preserved by flatness.  Thus
the small calculation determines the exponent count at generic quantum
parameters, for every smooth cubic threefold.

\subsection{Vanishing on centers through dimension two}

\begin{proposition}[Vanishing on low-dimensional centers]
\label{prop:atomic-lowdim}
\coverage{fragment}%
\lean{rankTwoResidueMarker_lowDimensionalOccurrenceNullity}%
\uses{prop:rank2-rigidity}%
\imports{BeauvilleSurfaces:minimal-surface-classification}%
For every smooth projective \(Z\) of dimension at most two,
\(I_{\mathrm{exp}}(Z)=I_{\mathrm{lat}}(Z)=0\).  The same vanishing holds
for every center copy in a quantum comparison.
\end{proposition}

\begin{proof}
It suffices to prove vanishing for \(I_{\mathrm{lat}}\), since
\(I_{\mathrm{exp}}\leq I_{\mathrm{lat}}\).
\emph{Points and curves.}
A point has one rank-one block.  For \(\PP^1\), the relation
\(p\star p=Qe^{t_p}\) gives two rank-one generic blocks.

Let \(C\) have genus \(g\ge1\).  There is no nonconstant genus-zero map to
\(C\), so its even big quantum product is classical.  With
\(\int_Cp=1\) and \(a=2-2g\), its centered even connection is
\[
 z\partial_zY=
 \left[z^{-1}aE_{21}+\diag(1/2,-1/2)\right]Y.
\]
For \(g=1\), \(N=0\).  For \(g>1\), the modification
\(S=\diag(z,1)\) gives residue \(aE_{21}-\tfrac12I\), of discriminant zero.
Its two eigenvalues, and therefore its two exponent classes, are equal.  Thus
every curve has zero invariant.

\smallskip
\noindent\emph{Surfaces with nef canonical class.}
Let \(S\) be a surface with \(K_S\) nef.  For a quantum-product term with
input degree \(b\), non-scalar Euler degree \(a\ge2\), curve class \(d\), and
nonunit even bulk degrees \(e_\nu\ge2\), the dimension axiom gives output
degree
\[
 a+b-2c_1(S)\cdot d+\sum_\nu(e_\nu-2)\ge b+2.
\]
The centered Euler operator strictly raises ordinary degree.  Its only block
is the whole even cohomology, of rank \(b_0+b_2+b_4\ge3\), so its invariant is
zero.

\smallskip
\noindent\emph{Other surfaces.}
A minimal surface with non-nef canonical class is \(\PP^2\) or a ruled
surface \(\PP_C(V)\) \cite[Chapter~VI]{BeauvilleSurfaces}.  The relation
\(H^{\star3}=Q\) makes \(\PP^2\) generically semisimple.  Formula
\eqref{eq:intrinsic-count-formulas} and the curve case give zero for ruled surfaces.
Every nonminimal surface is obtained from a minimal one by point blowups;
formula~\eqref{eq:intrinsic-count-formulas} adds only zero point terms.  This proves
intrinsic vanishing for every surface.
Lemma~\ref{lem:faithful-center-base-change} identifies every separate center
summand with a faithful scalar extension of the corresponding intrinsic generic
block, and the invariant is unchanged by that extension.
\end{proof}

\subsection{The projective endpoint and the contradiction}

Formula~\eqref{eq:intrinsic-count-formulas} applied to
\(X\times\PP^1=\PP_X(\OO_X^{\oplus2})\) gives
\begin{equation}
 I_{\mathrm{exp}}(X\times\PP^1)=2.
 \label{eq:atomic-product-value}
\end{equation}
On the other hand, over the generic numerical Novikov field
\(\Lambda=\operatorname{Frac}\C[[Q]]\),
\(QH^\bullet(\PP^4)=\Lambda[H]/(H^5-Q)\).  At \(Q\ne0\), multiplication
by \(5H\) has five distinct eigenvalues after algebraic closure.  Thus
\begin{equation}
 I_{\mathrm{exp}}(\PP^4)=0.
 \label{eq:atomic-projective-value}
\end{equation}

\begin{proof}[Proof of Theorem~\ref{thm:every-cubic}]
Proposition~\ref{prop:atomic-lowdim} makes every actual center correction zero
in a weak factorization of smooth projective fourfolds.  Hence
projective weak factorization and
\eqref{eq:intrinsic-count-formulas} make \(I_{\mathrm{exp}}\) birationally
invariant.  If \(X\times\PP^1\) were rational, then
\[
 2=I_{\mathrm{exp}}(X\times\PP^1)
 =I_{\mathrm{exp}}(\PP^4)=0,
\]
a contradiction.
\proves{thm:every-cubic}%
\end{proof}
\subsection{The additive formulation}
For retained block data \(\mathcal T\), let \(\mathcal L_{\mathcal T}\)
be the monoid of finite multisets of regular-isomorphism classes of blocks,
with addition by multiset union.  Write \(\mathcal B_{\mathcal T}(Y)\) for the
block multiset of \(Y\).  An \emph{additive invariant} is a homomorphism
\[
 w:\mathcal L_{\mathcal T}\longrightarrow A,
 \qquad I_w(Y):=w(\mathcal B_{\mathcal T}(Y)),
\]
where \(A\) is a commutative monoid, possibly without subtraction.  For a blowup
\(\beta:\Bl_ZY\to Y\) of codimension \(c\ge2\), retain the data of each
center occurrence separately:
\[
 \omega_j=(\beta,Z,j,\varsigma_j,\chi_j),
 \qquad 0\le j\le c-2.
\]
Here \(\varsigma_j\) is the reconstruction parameter and \(\chi_j\) its
numerical Novikov specialization.  Set
\(c_w(\omega_j):=w(\mathcal B_{\mathcal T}(Z;\omega_j))\);
equality of these values need not identify the block multisets.

\begin{theorem}[A birational invariance criterion]
\label{thm:marker-ledger}
\coverage{conditional_deduction}%
\lean{effectiveBlockLedger_fold_unique,%
  occurrenceIndexedMarker_eq_of_birational,%
  occurrenceIndexedMarker_descends_to_birationalClasses}%
\imports{AKMW:weak-factorization}%
Let \(w:\mathcal L_{\mathcal T}\to A\) be an additive invariant unchanged by the
allowed regular scalar extensions and formal coordinate changes.  Suppose
\begin{align}
 I_w(\PP_Y(V))&=rI_w(Y),
 &&\rk V=r,                                      \label{eq:marker-pbundle}\\
 I_w(\Bl_ZY)&=I_w(Y)+\sum_{j=0}^{c-2}c_w(\omega_j),
 &&\operatorname{codim}_YZ=c\ge2.                \label{eq:marker-blowup}
\end{align}
If every summand associated with a smooth center of dimension at most \(d-2\)
has zero correction, then \(I_w\) is a birational invariant of smooth
projective \(d\)-folds.
\end{theorem}

\begin{proof}
Projective weak factorization writes a birational map of smooth projective
\(d\)-folds as a zigzag of blowups and blowdowns along smooth centers of
codimension at least two \cite[Theorem~0.3.1]{AKMW}.  By
\eqref{eq:marker-blowup} and the vanishing of every center correction, the
invariant is unchanged at each step; the equalities therefore telescope.
Disconnected centers are treated componentwise, while a divisor blowup is an
isomorphism.
\end{proof}

\begin{proposition}[Marker formulas for the QDM comparisons]
\label{prop:qdm-operation-ledgers}
\coverage{conditional_deduction}%
\lean{qdmComparison_evenCarrierEquiv,%
  qdmProjectiveBundle_markerFormula_of_suspensionComparison,%
  qdmBlowup_markerFormula_of_suspensionComparison,%
  qdmProjectiveBundle_markerFormula_of_blockComparison,%
  qdmBlowup_markerFormula_of_occurrenceBlockComparison,%
  qdmBlowup_step_of_occurrenceBlockComparison}%
\imports{IritaniKoto:thm-5.1,IritaniBlowup:thm-5.18}%
Let an additive invariant be preserved by the regular scalar extensions and formal
coordinate changes described above, and by every recorded grading suspension
if its block type retains grading.  The projective-bundle and blowup QDM
comparisons on even cohomology then supply \eqref{eq:marker-pbundle} and
\eqref{eq:marker-blowup}; the contribution of the \(j\)-th center summand is
evaluated using its separate data \(\omega_j\).
\end{proposition}

\begin{proof}
The comparison proof of Proposition~\ref{prop:intrinsic-count-formulas}
uses only regular transport, generic separation and additivity.
Apply it to \(w\), retaining each center occurrence and any grading
suspensions in its block data.
\uses{prop:intrinsic-count-formulas}%
\imports{IritaniKoto:thm-5.1,IritaniBlowup:thm-5.18}%
\end{proof}
